Semantic fraud has become a practical risk in mobile banking transfer scenarios. Conventional risk-control systems usually rely on structured indicators such as transaction amount, device status, login location, and transaction frequency. These indicators are useful for identifying many abnormal operations, but they are less effective when the user is manipulated by social-engineering tactics and completes the transfer through his or her own account and device. In such cases, the transaction parameters may appear normal, while the actual risk is reflected in the transfer intention, the user's supplementary explanation, and subtle behavioral changes during the operation. To address this problem, this thesis designs and implements a prototype semantic fraud risk-control system for mobile banking based on AI Agent and the Model Context Protocol (MCP). The system integrates semantic analysis, micro-behavior perception, rule-based governance, and audit tracing into a unified transfer risk-control process.

The proposed system adopts a layered HIRD architecture, which consists of the harvesting integration layer, the intelligence decision layer, the regulation governance layer, and the delivery execution layer. The harvesting layer collects micro-behavioral features during the transfer process, including input duration, unexpected pauses, content modifications, page dwell time, and operation rhythm. These features are combined with transaction amount, payee information, device metadata, and transaction remarks to form a risk snapshot. The intelligence layer uses AI Agent to analyze transaction texts and user explanations. Retrieval-Augmented Generation (RAG) is introduced to provide relevant fraud cases and rule fragments, so that the analysis does not depend only on the internal knowledge of the large language model. The regulation layer does not directly use open-ended model responses as business decisions. Instead, it combines semantic risk scores, structured rule scores, and scenario adjustment factors to produce interpretable risk results, and then routes each transaction to pass, secondary verification, or block. The delivery layer manages transaction state transitions, manual review, audit logging, and sensitive information masking, ensuring that the Agent is responsible for analysis rather than direct fund-related execution.

In the implementation, FastAPI is used to build the asynchronous backend service and to separate the hot pre-check path from the slower semantic-analysis path. The mobile prototype demonstrates transfer submission, risk warning, secondary verification, and review feedback. MCP is used to encapsulate atomic tools such as risk calculation, audit logging, and business execution, providing standardized input-output formats, idempotency checks, and Trace ID based tracking. The system is evaluated with 66 backend automated test cases and 54 core scenario samples. The results show that the prototype achieves an exact-match rate of 92.6\%, a fraud recall rate of 100\%, and a direct false-blocking rate of 0\% for normal transactions on the constructed scenario set. In the latency test, the average response time of the pre-check hot path is 43.60 ms, with a P95 latency of 78.39 ms. The average response time of the secondary verification path is 2360.31 ms, with a P95 latency of 2767.42 ms. These results indicate that the prototype can complete semantic fraud identification, graded intervention, and audit tracing under the current test conditions. At the same time, the evaluation is still limited by the scale of constructed samples, the lack of real banking transaction data, and the stability of external model calls, which should be further examined in future work.
